\documentclass[10pt,a4paper]{amsart}
\usepackage[margin=19mm]{geometry}
\usepackage{amsmath,amssymb,mathtools}
\usepackage[scr=rsfs,cal=euler]{mathalfa}
\usepackage{xfrac}
\usepackage[unicode,hidelinks,bookmarks=false]{hyperref}
\newcommand{\ie}{i.e.\ }
\newcommand{\cf}{cf.\ }
\newcommand{\resp}{respectively\ }
\newcommand{\Subsection}{\subsection}
\providecommand{\nobreakdash}{\nobreak\hbox{-}}
\setlength{\emergencystretch}{2em}
\multlinegap0pt
\usepackage{amssymb}
\newtheorem{lemma}{Lemma}
\title{Qualitative Aspects of Periodic Traveling Waves for the Sinh-Gordon equation}
\author{Beatriz Signori Lonardoni, Fabio Natali}
\date{Source: arXiv:2601.01923v1 (CC BY 4.0)}
\begin{document}
\maketitle

\noindent\emph{Source and licence.} Qualitative Aspects of Periodic Traveling Waves for the Sinh-Gordon equation. Version arXiv:2601.01923v1 (CC BY 4.0), distributed by arXiv under the Creative Commons Attribution 4.0 International licence, http://creativecommons.org/licenses/by/4.0/, as recorded in the arXiv licence metadata for this exact version and shown on the abstract page https://arxiv.org/abs/2601.01923v1. Full licence notice: \texttt{licenses/arxiv-2601.01923-CC-BY-4.0.txt}.\par\smallskip

\noindent\textit{Editorial scope.} Counted unit: the article's lemma lem.montain (source0.tex lines 686--688). The article's complete original proof of it is delivered with it (source0.tex lines 689--741, one \texttt{proof} environment of 53 lines). No mathematics is written, repaired or re-derived: the statement and the proof are the article's own lines, in the article's own notation, and no retained line is substituted or re-flowed.  Retained uncounted as the source-local context the proof needs: lines 616--618.  Nothing was omitted from any retained span, so the package carries no omission record.  No retained line contains a citation, so the wrapper carries no bibliography and no bibliography entry.  No retained line contains a figure environment, an image-inclusion command or drawing code, so no image is delivered and the package declares no asset dependency.  Editorial formatting only, in addition: an AMS \texttt{amsart} preamble that restates the article's own declarations and macros used by the retained lines, namely the environments \texttt{lemma} on separate counters, together with the standard packages those lines need.  The retained environments print in this document as Lemma 1 in order of appearance, whereas the article prints the same items with the numbering of its own full document.  The complete native source of the article as distributed in the arXiv e-print for this exact version is bundled unchanged as \texttt{sources/192120/source0.tex}.
 \begin{equation}\label{functionalS}
 	\mathcal{S}_\mathtt{w}(v)= \int_{0}^{L}\left[ \frac{(v')^2}{2}-\frac{1}{\mathtt{w}} \cosh(v) +\frac{1}{\mathtt{w}}\right]dx.
 \end{equation}

 \begin{lemma}\label{lem.montain}
Let $\mathtt{w}\in(0,1)$ be fixed and  consider  $L\in(0,2\pi \sqrt{\mathtt{w}})$. The functional $\mathcal{S}_\mathtt{w}$ defined in \eqref{functionalS} has the mountain pass geometry. 
 \end{lemma}

 \begin{proof}
 	Indeed, using the Taylor formula with the Lagrange remainder to the hyperbolic cosine function around zero, we obtain a constant $\mathtt{M}(L)>0$ depending on $L$ satisfying %$|\partial_x^4\cosh(x) w^4 |\leq \mathtt{M}(L)|w|^4 $	 for all $x\in(0,L)$	
 	\begin{equation}\label{lem.montain.eq2}
 		\cosh(v)-\cosh(0)=\sinh(0) v+\frac{1}{2!} \cosh(0)v^2+\frac{1}{3!}\sinh(0)v^3+r_3(v),
 	\end{equation}where 
 	\begin{equation}\label{lem.montain.eq2'}
 		|r_3(v)|\leq \frac{\mathtt{M}(L)}{4!}|v|^4. 
 	\end{equation}
 	By the continuous embedding of $H^1_{per,m}\hookrightarrow L^4_{per.m}$, we find a constant $\mathtt{M}_0(L)>0$ such that 
 	\begin{equation}\label{lem.montain.eq3}
 		\|v\|_{L^4_{per}}\leq\mathtt{M}_0(L)  	\|v\|_{H^1_{per}}. 
 	\end{equation}  	 	
 	In addition, it follows from the Poincaré–Wirtinger inequality in $H_{per,m}^1$ that
 	\begin{align}  \label{PWineq}
 		\|v\|^2_{L^2_{per}}\leq 	 \biggl(\frac{L}{2\pi}\biggl)^2\|v\|^2_{H^1_{per}}.
 	\end{align} 
 	Substituting \eqref{lem.montain.eq2}-\eqref{PWineq} into \eqref{functionalS}, we obtain, for every $v\in H^1_{per,m}$,
 	\begin{equation}\label{lem.montain.eq4}\begin{array}{llllllll}
 		\mathcal{S}_\mathtt{w}(v)&=&\displaystyle \frac{1}{2} \|v\|^2_{H^1_{per}} -\frac{1}{\mathtt{w}}\int_{0}^{L} [\cosh(v) -1] dx  \\\\ &= & \displaystyle \frac{1}{2} \|v\|^2_{H^1_{per}} -\frac{1}{\mathtt{w}}\int_{0}^{L} \left|\frac{v^2}{2}+r_3(v)\right| dx \\\\ 
 		&\ge &\displaystyle  \frac{1}{2} \|v\|^2_{H^1_{per}} -\frac{1}{2\mathtt{w}}\|v\|^2_{L^2_{per}}  - \frac{\mathtt{M}(L)}{24\mathtt{w}}\|v\|^4_{L^4_{per}} \\\\&\ge &\displaystyle  \frac{1}{2} \|v\|^2_{H^1_{per}} -\frac{1}{2\mathtt{w}}	 \biggl(\frac{L}{2\pi}\biggl)^2\|v\|^2_{H^1_{per}}  - \frac{\mathtt{M}(L)\mathtt{M}_0^4(L)}{24\mathtt{w}}\|v\|^4_{H^1_{per}} \\\\&\ge & \displaystyle \frac{1}{2}\left[1 - \frac{1}{\mathtt{w}}\biggl(\frac{L}{2\pi}\biggl)^2	 \right]\|v\|^2_{H^1_{per}}  - \frac{\mathtt{M}(L)\mathtt{M}_0^4(L)}{24\mathtt{w}}\|v\|^4_{H^1_{per}}.
 	\end{array}\end{equation}
 	Thus, since $L\in(0,2\pi \sqrt{\mathtt{w}})$, we can assert the existence of $\varepsilon>0$ small enough such that
 	\begin{equation}
 		0<	\varepsilon<\sqrt{\frac{12\mathtt{w}}{\mathtt{M}(L)\mathtt{M}_0^4(L)}\left[1 - \frac{1}{\mathtt{w}}\biggl(\frac{L}{2\pi}\biggl)^2	 \right]}.
 	\end{equation}
 	\indent Therefore, for any $ v\in H^1_{per,m}$ satisfying $\|v\| _{H^1_{per}}=\varepsilon$, we have from \eqref{lem.montain.eq4} that $	\mathcal{S}_\mathtt{w}(v) > 0$. This means that 
 	\begin{equation}\label{lem.montain.eq5}
 		\displaystyle	\inf_{\|v\|_{H^1_{per}}=\varepsilon} \mathcal{S}_\mathtt{w}(v)>0=\mathcal{S}_\mathtt{w}(0).
 	\end{equation}
 	In addition, for any $v\in H^1_{per,m}$ and $\tau>0$, 
 	\begin{equation}\label{lem.montain.eq1}
 		\mathcal{S}_\mathtt{w} (\tau v)=\frac{\tau^2}{2} \|v\|^2_{H^1_{per}}+\frac{L}{\mathtt{w}}-\frac{1}{\mathtt{w}}\int_{0}^{L}\cosh(\tau v)dx.
 	\end{equation}
 	Since the hyperbolic cosine function is positive and exhibits exponential-type growth, it follows from \eqref{lem.montain.eq1} that $	\mathcal{S}_\mathtt{w} (\tau v)<0$ if $\tau$ is large enough. Let $\tau_0>0$ be such that  $	\mathcal{S}_\mathtt{w} (\tau_0 v)<0$ and $\|\tau_0v\|_{H^1_{per}}>\varepsilon$. By defining $e=(\tau_0 v) \in H^1_{per,m}$ and $\displaystyle r_e=\varepsilon$, we prove that	   	 	\begin{equation}\label{lem.montain.eq1'} \inf_{\|v\|=r_e}\mathcal{S}_\mathtt{w}(v)>\mathcal{S}_\mathtt{w}(0)\ge \mathcal{S}_\mathtt{w}(e). 
 	\end{equation}
 	\indent Now, let us consider the mapping
 	$\gamma:[0,1]\to H^1_{per,m}$ defined as $\gamma(\tau)=\tau_0 \tau v$. It follows that $\gamma\in C([0,1], H^1_{per,m})$, $\gamma(0)=0$ and $\gamma(1)=e$, proving that $\Gamma$ is nonempty. 	   
 	This implies that for any $\gamma\in \Gamma$, we have $\|\gamma(1)\|_{H^1_{per}} >r_e>0=\|\gamma(0)\|_{H^1_{per}}$. By continuity of $\gamma$, there exists $\tau _\gamma\in (0,1)$ such that $\|\gamma(\tau_ \gamma)\|_{H^1_{per}} =r_e$. Therefore, from \eqref{lem.montain.eq5} we obtain
 	\begin{equation}
 		\max_{\tau\in [0,1]}\mathcal{S}_\mathtt{w}(\gamma(\tau))\ge\mathcal{S}_\mathtt{w}(\gamma(\tau_\gamma))>0, \quad \gamma\in \Gamma,
 	\end{equation}consequently, the mountain pass level satisfies
 	\begin{align}\label{level}
 		\tilde{c}=\inf_{\gamma\in\Gamma}\max_{\tau\in[0,1]} \mathcal{S}_\mathtt{w}(\gamma(\tau))>0.
 	\end{align}
 	Moreover,
 	\begin{equation}\label{lem.montain.eq6} \inf_{\|v\|=r_e}\mathcal{S}_\mathtt{w}(v)\leq\mathcal{S}_\mathtt{w}(\gamma(\tau_\gamma))\leq \max_{\tau\in[0,1]} \mathcal{S}_\mathtt{w}(\gamma(\tau)), \quad\gamma\in \Gamma,\nonumber
 	\end{equation}and by the definition of the infimum, we obtain
 	\begin{equation}\label{lem.montain.eq7} 0<\inf_{\|v\|=r_e}\mathcal{S}_\mathtt{w}(v) \leq \inf_{\gamma\in\Gamma}\max_{\tau\in[0,1]} \mathcal{S}_\mathtt{w}(\gamma(\tau)) =\tilde{c}.
 	\end{equation}In particular, considering the path $\gamma_0(\tau)=\tau e \in \Gamma$ 		  
 	in \eqref{lem.montain.eq7} and using \eqref{lem.montain.eq1}, it follows that
 	\begin{equation}\label{lem.montain.eq8} 0<\inf_{\|v\|=r_e}\mathcal{S}_\mathtt{w}(v) \leq\tilde{c}\leq     \max_{\tau\in[0,1]}\mathcal{S}_\mathtt{w}(\gamma_0(\tau)) =\max_{\tau\in[0,1]}\mathcal{S}_\mathtt{w}(\tau e)\ \leq \frac{1}{2}\|e\|^2_{H^1_{per}}+\frac{L}{\mathtt{w}}<+\infty,
 	\end{equation}
 	 that is, $\mathcal{S}_\mathtt{w}$ has a mountain pass geometry. \end{proof}

\end{document}
